\documentclass[a4paper,11pt]{ltjsarticle}
\usepackage{luatexja}
\usepackage{amsmath,amssymb,amsthm}
\usepackage{longtable,booktabs,array}
\usepackage{hyperref}
\usepackage[margin=25mm]{geometry}

\newtheorem{theorem}{定理}[section]
\newtheorem{proposition}[theorem]{命題}
\newtheorem{lemma}[theorem]{補題}
\newtheorem{corollary}[theorem]{系}
\newtheorem{definition}[theorem]{定義}
\newtheorem{remark}[theorem]{注意}
\newtheorem{observation}[theorem]{観察}

\providecommand{\tightlist}{\setlength{\itemsep}{0pt}\setlength{\parskip}{0pt}}
\newcounter{none}


\begin{document}


\section{D1 v2.0:
中心投影フレームワークの粒子データ構造仕様------周波数解釈と動力学拡張}\label{d1-v2.0-ux4e2dux5fc3ux6295ux5f71ux30d5ux30ecux30fcux30e0ux30efux30fcux30afux306eux7c92ux5b50ux30c7ux30fcux30bfux69cbux9020ux4ed5ux69d8ux5468ux6ce2ux6570ux89e3ux91c8ux3068ux52d5ux529bux5b66ux62e1ux5f35}

\textbf{著者}: 木原 範昭 (Noriaki Kihara) \textbf{所属}: WF System Co.,
Ltd. \textbf{ORCID}: 0009-0004-6753-4020 \textbf{版}: v2.1
\textbf{日付}: 2026 年 5 月 \textbf{DOI}:
\href{https://doi.org/10.5281/zenodo.19941250}{10.5281/zenodo.19941250}（v2.0、v2.1
公開時に更新予定） \textbf{Concept DOI}:
\href{https://doi.org/10.5281/zenodo.19941249}{10.5281/zenodo.19941249}
\textbf{先行版}: v1.0 (2026年4月、未公開) / v2.0 (2026年5月公開)

\begin{center}\rule{0.5\linewidth}{0.5pt}\end{center}

\subsection{§0
本稿の位置づけ}\label{ux672cux7a3fux306eux4f4dux7f6eux3065ux3051}

本稿は数学的証明を含まない設計仕様である。先行論文群
{[}W1{]}--{[}W11{]}, {[}M1{]}--{[}M3{]}, {[}補講1{]}--{[}補講4{]},
{[}思考実験\_R軸{]}
において暗黙に前提とされている粒子のデータ構造を、型付き整数演算の設計仕様として明示的に定義する。

\subsubsection{v2.0
からの改訂（v2.1）}\label{v2.0-ux304bux3089ux306eux6539ux8a02v2.1}

\textbf{\texttt{δδt} フィールドの削除}：v2.0 で Particle 構造に追加した
\texttt{δδt\ :\ DeltaT}（処理周期変化、加速度相当）を\textbf{削除}する。理由は以下の通り：

\begin{enumerate}
\def\labelenumi{\arabic{enumi}.}
\tightlist
\item
  \textbf{可逆性の破れ}：\texttt{δt\ +=\ δδt} を含む \texttt{step(p)}
  は、δt の符号反転で元の状態に戻らない（{[}W3{]}
  の情報保存要件に反する）
\item
  \textbf{保存則との非整合}：単一粒子で δt
  が変化することは、外因なく固有時の進み方が変わることを意味し、エネルギー・運動量保存則の暗黙の破れを引き起こす
\item
  \textbf{GR
  自由粒子原理との整合}：自由粒子は測地線に従い、加速度効果は外力（相互作用）から生じる。単一粒子の基本演算で加速度を扱うべきではない
\end{enumerate}

加速度効果は \textbf{すべて相互作用処理（粒子間または粒子-Mesh
間）の中で、保存則を満たす形で σ や δt
が変化する}ことで表現される。詳細は別稿
D2「中心投影フレームワークの動力学規則仕様」（執筆中）を参照。

修正された Particle 構造は \textbf{4 フィールド}
\texttt{(σ,\ V,\ δV,\ δt)} となる（§5）。Mesh 構造からも対応する
\texttt{δδt\_mesh} を削除（§7）。

\subsubsection{v1.0 からの主要な改訂（v2.0 で導入、v2.1
で継続）}\label{v1.0-ux304bux3089ux306eux4e3bux8981ux306aux6539ux8a02v2.0-ux3067ux5c0eux5165v2.1-ux3067ux7d99ux7d9a}

\begin{enumerate}
\def\labelenumi{\arabic{enumi}.}
\tightlist
\item
  \textbf{σ.k\_xyzt
  の解釈を「波数・加速度」から「各軸の振動周波数」へ統一}（§1.2）
\item
  \textbf{Particle 構造に
  δV（位相変位）を追加}（§5）。〜\textasciitilde v2.0 では δδt
  も追加したが v2.1 で削除\textasciitilde{}
\item
  \textbf{Mesh 構造を新設し、縦波モード（ヒッグス・重力波）を Particle
  と分離}（§7）
\item
  \textbf{時間 t を特権化せず、6 軸対称性を完全に保持}（§3, §6）
\end{enumerate}

データ型システム（int256
ベース）と整数演算原則（除算・平方根・三角関数禁止）は v1.0
から変更なし。

\subsubsection{本稿の主張}\label{ux672cux7a3fux306eux4e3bux5f35}

\begin{enumerate}
\def\labelenumi{\arabic{enumi}.}
\tightlist
\item
  宇宙の有界性とプランクスケールの離散性から、全位相量が256ビット符号付き整数で表現可能であること（§1）
\item
  xyzt 軸の {[}L⁻¹{]}
  次元は「位相空間における各軸の振動周波数」として解釈するのが自然であり、6軸対称性を完全に保つこと（§3）
\item
  粒子の完全な状態が \texttt{SignVector}（個性ラベル）+
  \texttt{OneDWave{[}κ{]}}（位相状態）+ \texttt{δV}（位相変位）+
  \texttt{δt}（処理周期）の 4 フィールドで記述できること（§5）
\item
  ヒッグスを始めとする縦波モードは Particle ではなく \texttt{Mesh}
  の振動として記述すべきであり、これにより質量生成・重力波・宇宙論的振動を統一的に扱えること（§7）
\item
  全演算が整数の加算・乗算で閉じ、除算・平方根・三角関数を必要としないこと（§6）
\item
  基本演算 \texttt{step(p)}（V
  のみ更新）は時間反転対称・保存則整合・null 透過の 5
  原則すべてを満たすこと（§6）
\end{enumerate}

\subsubsection{本稿の主張範囲外}\label{ux672cux7a3fux306eux4e3bux5f35ux7bc4ux56f2ux5916}

\begin{itemize}
\tightlist
\item
  演算規則（位相発展の物理的意味、相互作用の具体形）の動力学的定式化は別稿（D2）に委ねる
\item
  具体的な物理量の数値的導出
\item
  ハードウェア実装の詳細
\item
  数学的定理の証明
\end{itemize}

\subsubsection{参照論文}\label{ux53c2ux7167ux8ad6ux6587}

\begin{itemize}
\tightlist
\item
  {[}W3{]} 万物の理論が満たすべき構造要件の定式化 (DOI:
  10.5281/zenodo.19601592)
\item
  {[}W8{]} 6次元超直方体の集合構造とその組合せ論的性質 (DOI:
  10.5281/zenodo.19762134)
\item
  {[}W9{]} sine-Gordon方程式 ── 位相的ソリトンの基礎理論 (DOI:
  10.5281/zenodo.19650966)
\item
  {[}W10{]} 形不変波の4つのモード (DOI: 10.5281/zenodo.19709798)
\item
  {[}W11{]} 形不変波の相互作用 (DOI: 10.5281/zenodo.19763463)
\item
  {[}M1{]} 中心投影の三層構造モデル: R--R₁--R₀ (DOI:
  10.5281/zenodo.19691713)
\item
  {[}思考実験\_R軸{]} 6次元符号化 xyztRQ の再検討 (DOI:
  10.5281/zenodo.19904714)
\end{itemize}

\begin{center}\rule{0.5\linewidth}{0.5pt}\end{center}

\subsection{§1 int256
有界性と次元解釈}\label{int256-ux6709ux754cux6027ux3068ux6b21ux5143ux89e3ux91c8}

\subsubsection{§1.1 有界性要件の根拠（v1.0
から変更なし）}\label{ux6709ux754cux6027ux8981ux4ef6ux306eux6839ux62e0v1.0-ux304bux3089ux5909ux66f4ux306aux3057}

{[}W3{]} は万物の理論の構造要件として有界性を要求する。{[}M1{]}
の三層構造 \(R\)--\(R_1\)--\(R_0\) において、最大の曲率半径 \(R_0\)
は有限である。この有限性とプランクスケールの離散性から、全位相量の表現に必要なビット数の上界が導かれる。

\subsubsection{§1.2
必要ビット数}\label{ux5fc5ux8981ux30d3ux30c3ux30c8ux6570}

\textbf{空間軸}: 観測可能宇宙径 \(D \approx 1.892 \times 10^{27}\)
m、プランク長 \(l_P = 1.616 \times 10^{-35}\) m から

\[n_{\text{space}} = D / l_P \approx 1.171 \times 10^{62}, \quad b_{\text{space}} = 207 \text{ bit}\]

\textbf{時間軸}: 宇宙時間 \(T \approx 8.710 \times 10^{17}\)
s、プランク時間 \(t_P = 5.391 \times 10^{-44}\) s から

\[n_{\text{time}} = T / t_P \approx 1.616 \times 10^{61}, \quad b_{\text{time}} = 204 \text{ bit}\]

\textbf{結論}: 256ビット符号付き整数で空間 14 桁・時間 15
桁の余裕を持って表現可能。

\subsubsection{§1.3
自然単位系と次元}\label{ux81eaux7136ux5358ux4f4dux7cfbux3068ux6b21ux5143}

\(c = 1\), \(G = 1\)（幾何学化単位系、\(\hbar\) は採用しない）の下で

\[[L] = [T] = [M], \quad [L^{-1}] = [M^{-1}]\]

\subsubsection{§1.4 周波数解釈（v2.0
の中核改訂）}\label{ux5468ux6ce2ux6570ux89e3ux91c8v2.0-ux306eux4e2dux6838ux6539ux8a02}

v1.0 では xyzt 軸の \([L^{-1}]\)
を「波数」「加速度」と混在的に解釈していた。v2.0
では位相空間の自然な解釈として \textbf{「各軸方向の振動周波数」}
に統一する。

{\def\LTcaptype{none} % do not increment counter
\begin{longtable}[]{@{}
  >{\raggedright\arraybackslash}p{(\linewidth - 6\tabcolsep) * \real{0.2500}}
  >{\raggedright\arraybackslash}p{(\linewidth - 6\tabcolsep) * \real{0.2500}}
  >{\raggedright\arraybackslash}p{(\linewidth - 6\tabcolsep) * \real{0.2500}}
  >{\raggedright\arraybackslash}p{(\linewidth - 6\tabcolsep) * \real{0.2500}}@{}}
\toprule\noalign{}
\begin{minipage}[b]{\linewidth}\raggedright
物理量
\end{minipage} & \begin{minipage}[b]{\linewidth}\raggedright
次元
\end{minipage} & \begin{minipage}[b]{\linewidth}\raggedright
v1.0 解釈
\end{minipage} & \begin{minipage}[b]{\linewidth}\raggedright
\textbf{v2.1 解釈}
\end{minipage} \\
\midrule\noalign{}
\endhead
\bottomrule\noalign{}
\endlastfoot
σ.k\_x, k\_y, k\_z & \([L^{-1}]\) & 空間波数・加速度 &
\textbf{各空間軸の振動周波数} \\
σ.k\_t & \([L^{-1}]\) & 時間振動数 & \textbf{t
軸の振動周波数}（他軸と対等） \\
σ.R & \([L]=[M]\) & 曲率半径 & \textbf{スケール量・対応周期}（双対） \\
σ.Q & \([1]\) & 内部対称性 & 内部対称性（変更なし） \\
δt & \([L]=[M]\) & 固有時間ステップ &
\textbf{処理ステップ周期}（時間ではない） \\
\textsubscript{δδt} & \textsubscript{\([L]=[M]\)} &
\textasciitilde（v1.0 になし）\textasciitilde{} & \textasciitilde v2.0
で追加、v2.1 で削除（保存則・可逆性に反するため）\textasciitilde{} \\
\end{longtable}
}

\subsubsection{§1.5
周波数解釈の優位性}\label{ux5468ux6ce2ux6570ux89e3ux91c8ux306eux512aux4f4dux6027}

\textbf{(1) 6軸対称性の完全保持}: 位置型 4
軸（xyzt）はそれぞれ独立な振動モードを持ち、t 軸も他軸と対等。

\textbf{(2) 連続体力学概念の不要化}:
「加速度」は位置の二階時間微分という連続体概念であり、離散整数演算の枠組みには馴染まない。「振動周波数」は位相空間の自然な離散量。

\textbf{(3) 加速度効果は相互作用として処理}:
加速度的な効果（重力場・慣性系の影響）は単一粒子では表現せず、\textbf{すべて相互作用処理}（粒子間または粒子-Mesh
間）の中で σ や δt の変化として保存則整合的に扱う。これは GR
の自由粒子原理（自由粒子は測地線に従う）と整合する。詳細は別稿 D2
を参照。

\textbf{(4) 周波数 × 周期 = 位相}: \(\sigma.k \cdot \delta t\)
が無次元の位相増分となり（\([L^{-1}] \cdot [L] = [1]\)）、位相発展則が型整合する。

\textbf{(5) 特殊相対論の自動内包}: dispersion relation
\(\sigma.k_t^2 = \sigma.k_x^2 + \sigma.k_y^2 + \sigma.k_z^2 + m^2\)
が周波数間の整数関係として自然に書ける。

\begin{center}\rule{0.5\linewidth}{0.5pt}\end{center}

\subsection{§2 メタ型の定義}\label{ux30e1ux30bfux578bux306eux5b9aux7fa9}

\subsubsection{§2.1 スカラー型（v1.0
から変更なし）}\label{ux30b9ux30abux30e9ux30fcux578bv1.0-ux304bux3089ux5909ux66f4ux306aux3057}

\begin{verbatim}
型名         次元       値域                      用途
L256         [L]=[M]   |v| ≤ 2^255 - 1          長さ・時間・質量・周期
Linv256      [L⁻¹]     |v| ≤ 2^255 - 1          周波数・スケール逆数
Scalar256    [1]       |v| ≤ 2^255 - 1          無次元量
FullPhase    [1]       0 ≤ v ≤ P                0–720°
CyclePhase   [1]       0 ≤ v ≤ P/2              0–360°（位置位相）
DeltaPhase   [1]       -P/2 ≤ v ≤ P/2           ±360°（位相変位）
SpinPhase    [1]       v ∈ {0, P/4, P/2, P}     {0, 180, 360, 720}°
\end{verbatim}

位相定数 \(P = \lfloor (2^{255}-1)/720 \rfloor \times 720\)。

\subsubsection{§2.2 軸成分型（v2.0
で再命名・解釈変更）}\label{ux8ef8ux6210ux5206ux578bv2.0-ux3067ux518dux547dux540dux89e3ux91c8ux5909ux66f4}

{\def\LTcaptype{none} % do not increment counter
\begin{longtable}[]{@{}
  >{\raggedright\arraybackslash}p{(\linewidth - 6\tabcolsep) * \real{0.2500}}
  >{\raggedright\arraybackslash}p{(\linewidth - 6\tabcolsep) * \real{0.2500}}
  >{\raggedright\arraybackslash}p{(\linewidth - 6\tabcolsep) * \real{0.2500}}
  >{\raggedright\arraybackslash}p{(\linewidth - 6\tabcolsep) * \real{0.2500}}@{}}
\toprule\noalign{}
\begin{minipage}[b]{\linewidth}\raggedright
型
\end{minipage} & \begin{minipage}[b]{\linewidth}\raggedright
基底
\end{minipage} & \begin{minipage}[b]{\linewidth}\raggedright
次元
\end{minipage} & \begin{minipage}[b]{\linewidth}\raggedright
物理的意味（v2.0）
\end{minipage} \\
\midrule\noalign{}
\endhead
\bottomrule\noalign{}
\endlastfoot
\texttt{SpatialFreq} \textless: Linv256 & int256 & \([L^{-1}]\) & x, y,
z 軸の振動周波数 \\
\texttt{TemporalFreq} \textless: Linv256 & int256 & \([L^{-1}]\) & t
軸の振動周波数（他軸と対等） \\
\texttt{CurvatureR} \textless: L256 & int256 & \([L]\) &
スケール量・周期（符号付き平方数） \\
\texttt{ColorQ} \textless: Scalar256 & int256 & \([1]\) &
内部対称ラベル（有限群） \\
\texttt{DeltaT} \textless: L256 & int256 & \([L]\) & 処理ステップ周期 \\
\end{longtable}
}

\textbf{v1.0 からの改名}: \texttt{SpatialK} →
\texttt{SpatialFreq}、\texttt{TemporalK} →
\texttt{TemporalFreq}。型システム自体は不変、解釈の明確化のみ。

\subsubsection{§2.3 演算規則（v1.0
から変更なし）}\label{ux6f14ux7b97ux898fux5247v1.0-ux304bux3089ux5909ux66f4ux306aux3057}

\textbf{加算}: 同型同士のみ許可（CyclePhase + DeltaPhase = CyclePhase
など）。

\textbf{乗算}: - \texttt{L256\ ×\ L256\ →\ L2\_256} \([L^2]\) -
\texttt{Linv256\ ×\ Linv256\ →\ Linv2\_256} \([L^{-2}]\) -
\texttt{Linv256\ ×\ L256\ →\ Scalar256} \([1]\)（無次元化）

\textbf{除算}: 使用しない。

\begin{center}\rule{0.5\linewidth}{0.5pt}\end{center}

\subsection{§3 SignVector
の定義（周波数解釈版）}\label{signvector-ux306eux5b9aux7fa9ux5468ux6ce2ux6570ux89e3ux91c8ux7248}

\subsubsection{§3.1 構造}\label{ux69cbux9020}

\begin{verbatim}
SignVector = (
    x : SpatialFreq,    -- x 軸の振動周波数 [L⁻¹]
    y : SpatialFreq,    -- y 軸の振動周波数
    z : SpatialFreq,    -- z 軸の振動周波数
    t : TemporalFreq,   -- t 軸の振動周波数 [L⁻¹]（他軸と対等）
    R : CurvatureR,     -- スケール量・周期 [L]
    Q : ColorQ,         -- 内部対称ラベル [1]
)
\end{verbatim}

\subsubsection{§3.2 物理的意味}\label{ux7269ux7406ux7684ux610fux5473}

各 σ 成分は「該当軸での振動周波数」または「スケール量」：

\begin{itemize}
\tightlist
\item
  σ.k\_axis = 0: 該当軸は振動しない（その方向の構造を持たない）
\item
  σ.k\_axis ≠ 0: 該当軸が周波数 \textbar k\_axis\textbar{} で振動する
\item
  符号 sign(k\_axis): 振動の位相方向（+ なら時計回り、− なら反時計回り）
\item
  \textbar k\_axis\textbar: 振動周波数の大きさ（離散整数値）
\end{itemize}

\subsubsection{§3.3 不変量}\label{ux4e0dux5909ux91cf}

SignVector
の値は相互作用イベントで変化するが、特定の組み合わせは安定モードを定義する：

\begin{itemize}
\tightlist
\item
  \(\kappa_s = |\{i \in \{x,y,z\} : k_i \neq 0\}|\) --- 空間非ゼロ数
\item
  \(\kappa_t = (1 \text{ if } k_t \neq 0 \text{ else } 0)\) ---
  時間非ゼロ指標
\item
  \(\kappa = \kappa_s + \kappa_t \in \{0, 1, 2\}\) --- {[}W10{]}
  命題2.1: \(\kappa \geq 3\) は不安定
\end{itemize}

\subsubsection{§3.4 dispersion relation}\label{dispersion-relation}

質量殻条件は周波数間の整数関係として直接表現：

\[\sigma.k_t^2 = \sigma.k_x^2 + \sigma.k_y^2 + \sigma.k_z^2 + m^2\]

\begin{itemize}
\tightlist
\item
  \(m = 0\)（光子・グルーオン）: \(k_t^2 = k_x^2 + k_y^2 + k_z^2\)
\item
  \(m \neq 0\)（質量粒子）: 差分の平方根が質量に対応
\end{itemize}

これにより特殊相対論が自動的に枠組みに含まれる。

\subsubsection{§3.5 W8 表 A との対応（v1.0
から継続）}\label{w8-ux8868-a-ux3068ux306eux5bfeux5fdcv1.0-ux304bux3089ux7d99ux7d9a}

{[}W8{]} 表 A の符号パターン \(\{+, -, 0\}^6\) が SignVector
の符号と対応する。整数値（軸スケール値）は {[}質量構造{]} §1 の \(v_i\)
に対応。

代表粒子：

{\def\LTcaptype{none} % do not increment counter
\begin{longtable}[]{@{}llll@{}}
\toprule\noalign{}
粒子 & 集合 \((x,y,z,t,R,Q)\) & \(\kappa\) & 分類 \\
\midrule\noalign{}
\endhead
\bottomrule\noalign{}
\endlastfoot
光子 \(\gamma\) & \((0,0,0,0,+,0)\) & 0 &
ゲージボソン（自由・球状波） \\
グルーオン \(g\) & \((0,0,0,0,0,Q_i \to Q_j)\) & 0 & ゲージボソン \\
ヒッグス \(H\) & \((0,0,+,0,0,0)\) & 1 & スカラー（縦波解釈は §7） \\
\(W^+\) & \((+,0,0,0,+,0)\) & 1 & ゲージボソン \\
\(Z^0\) & \((0,+,0,0,+,0)\) & 1 & ゲージボソン \\
重力子 & \((0,0,0,\pm,\pm,0)\) & 1 & テンソルボソン \\
電子 \(e^-\) & \((+,0,0,+,0,0)\) & 2 & フェルミオン \\
\end{longtable}
}

\begin{center}\rule{0.5\linewidth}{0.5pt}\end{center}

\subsection{§4 OneDWave の定義}\label{onedwave-ux306eux5b9aux7fa9}

\subsubsection{§4.1 構造（v1.0
から変更なし）}\label{ux69cbux9020v1.0-ux304bux3089ux5909ux66f4ux306aux3057}

\begin{verbatim}
OneDWave = (
    θ₀ : CyclePhase,   -- 位置位相 [1]      (S1 §2.1, W9 §8.2)
    A  : Scalar256,     -- 振幅 [1]          (S1 §2.1, W9 §8.2)
    ℓ  : DeltaPhase,    -- 広がり位相 [1]    (W9 §8.3)
)
\end{verbatim}

\subsubsection{§4.2 物理的意味}\label{ux7269ux7406ux7684ux610fux5473-1}

各粒子は非零位置型軸ごとに 1 つの OneDWave を持つ（最大 \(\kappa\)
個）。各 OneDWave はその軸方向の振動の現在位相状態を保持する。

\begin{itemize}
\tightlist
\item
  \(\theta_0\): 振動の現在位相（自由パラメータ、時間発展で変化）
\item
  \(A\): 振幅（自由パラメータ）
\item
  \(\ell\): 振動の局在広がり（SignVector から従属、{[}W9{]} §8.3）
\end{itemize}

\begin{center}\rule{0.5\linewidth}{0.5pt}\end{center}

\subsection{§5 Particle の定義（v2.1：4
フィールド構造）}\label{particle-ux306eux5b9aux7fa9v2.14-ux30d5ux30a3ux30fcux30ebux30c9ux69cbux9020}

\subsubsection{§5.1 構造}\label{ux69cbux9020-1}

\begin{verbatim}
Particle = (
    σ   : SignVector,    -- 個性（各軸の振動周波数）
    V   : OneDWave[κ],   -- 各非零位置型軸の現在位相状態
    δV  : OneDWave[κ],   -- 1 ステップでの位相変位（v2.0 で追加）
    δt  : DeltaT,         -- 現在の処理ステップ周期
)
\end{verbatim}

\textbf{v2.0 からの変更}: \texttt{δδt\ :\ DeltaT} フィールドを削除（5
フィールド → 4 フィールド）。理由は §0 の改訂説明を参照。

\subsubsection{§5.2
各フィールドの物理的意味}\label{ux5404ux30d5ux30a3ux30fcux30ebux30c9ux306eux7269ux7406ux7684ux610fux5473}

\paragraph{\texorpdfstring{\texttt{σ\ :\ SignVector}}{σ : SignVector}}\label{ux3c3-signvector}

粒子の個性。各軸でどの周波数で振動するか。相互作用イベントで値が変化する（消滅＋生成として記述、しかしデータ構造としては
σ の値が更新される）。

\paragraph{\texorpdfstring{\texttt{V\ :\ OneDWave{[}κ{]}}}{V : OneDWave{[}κ{]}}}\label{v-onedwaveux3ba}

各非零位置型軸の現在位相状態。配列長 \(\kappa\) は SignVector
の非零位置型軸数（\(\leq 2\)）。各 V{[}j{]} は対応軸の振動の (位置位相,
振幅, 広がり) を保持する。

インデックス対応：

\begin{verbatim}
σ.k_x ≠ 0 のとき V に wave_x が含まれる
σ.k_y ≠ 0 のとき V に wave_y が含まれる
σ.k_z ≠ 0 のとき V に wave_z が含まれる
σ.k_t ≠ 0 のとき V に wave_t が含まれる
\end{verbatim}

\paragraph{\texorpdfstring{\texttt{δV\ :\ OneDWave{[}κ{]}} (v2.0
で追加)}{δV : OneDWave{[}κ{]} (v2.0 で追加)}}\label{ux3b4v-onedwaveux3ba-v2.0-ux3067ux8ffdux52a0}

V と同型の差分構造。1 処理ステップで V がどう変化するかを保持する：

\begin{verbatim}
δV[j] = (
    Δθ₀ : DeltaPhase,    -- 位相変位
    ΔA  : Scalar256,      -- 振幅変位
    Δℓ  : DeltaPhase,     -- 広がり変位
)
\end{verbatim}

自由粒子の場合、δV{[}j{]}.Δθ₀ は σ から計算可能：

\[\delta V[j].\Delta\theta_0 = \sigma.k_{a(j)} \cdot \delta t\]

ただし相互作用イベントでは δV が独立に変更されうる。

\paragraph{\texorpdfstring{\texttt{δt\ :\ DeltaT}}{δt : DeltaT}}\label{ux3b4t-deltat}

\textbf{処理ステップ周期}（時間ではない）。普遍的な演算カウンタ。

\begin{itemize}
\tightlist
\item
  \(\delta t > 0\): 順方向演算（因果の順方向）
\item
  \(\delta t < 0\): 逆向き演算（可逆性の保証）
\item
  \(\delta t = 0\): 静止
\end{itemize}

時間として観測されるのは V{[}t-軸{]}.θ₀ の累積であり、δt
自身ではない（§6.4 参照）。

\paragraph{\texorpdfstring{\textasciitilde{}\texttt{δδt\ :\ DeltaT}（v2.1
で削除）\textasciitilde{}}{\textasciitilde δδt : DeltaT（v2.1 で削除）\textasciitilde{}}}\label{ux3b4ux3b4t-deltatv2.1-ux3067ux524aux9664}

v2.0 では Particle 構造に \texttt{δδt} フィールドを含めていたが、v2.1
で削除した。理由：

\begin{itemize}
\tightlist
\item
  \textbf{可逆性違反}: \texttt{step(p)} 内で \texttt{δt\ +=\ δδt}
  を実行すると、δt の符号反転による逆演算で元の状態に戻らない
\item
  \textbf{保存則違反}: 単一粒子で δt
  が変化することは外因なきエネルギー・運動量変化を意味し、保存則と整合しない
\item
  \textbf{GR 自由粒子原理}:
  加速度効果は外力（相互作用）の結果であるべきで、自由粒子の基本演算では扱えない
\end{itemize}

加速度効果は §3 以降の相互作用処理（粒子間または粒子-Mesh
間）の中で、保存則を満たす形で σ や δt
が変化することで表現される。詳細は別稿 D2 を参照。

\subsubsection{§5.3 自由粒子の動力学は δV
で完結}\label{ux81eaux7531ux7c92ux5b50ux306eux52d5ux529bux5b66ux306f-ux3b4v-ux3067ux5b8cux7d50}

物理法則が 2
階微分で閉じる（ニュートン力学・GR・量子力学）ことから当初は δδt
が必要と考えたが、v2.1 では：

\begin{itemize}
\tightlist
\item
  自由粒子は等速測地線運動のみ（GR 原理）
\item
  加速度は相互作用の結果
\item
  単一粒子の基本演算は \textbf{位相速度（δV）と現在位相（V）の更新のみ}
\end{itemize}

したがって \textbf{4 フィールド（σ, V, δV,
δt）で自由粒子表現として完全}。

\subsubsection{§5.4
メモリフットプリント（v2.1）}\label{ux30e1ux30e2ux30eaux30d5ux30c3ux30c8ux30d7ux30eaux30f3ux30c8v2.1}

\begin{verbatim}
σ      : 6 × 256 = 1,536 bit
V      : κ × 3 × 256 ≤ 1,536 bit (κ ≤ 2)
δV     : κ × 3 × 256 ≤ 1,536 bit
δt     : 256 bit
合計   ≤ 4,864 bit = 608 byte / 粒子
\end{verbatim}

v2.0 の 640 byte から 32 byte 削減（δδt フィールド削除分）。v1.0 の
224--416 byte
からは増加するが、依然として宇宙の任意の粒子を完全に表現する量として極めて少ない。

\begin{center}\rule{0.5\linewidth}{0.5pt}\end{center}

\subsection{§6 演算則}\label{ux6f14ux7b97ux5247}

\subsubsection{§6.1 基本：1
処理ステップでの更新（v2.1）}\label{ux57faux672c1-ux51e6ux7406ux30b9ux30c6ux30c3ux30d7ux3067ux306eux66f4ux65b0v2.1}

基本演算 \texttt{step(p)} は次の 1 つの操作のみ：

\begin{verbatim}
function step(p : Particle | null) → Particle | null:
    if p == null:
        return null
    for j in range(p.κ):
        p.V[j].θ₀ += p.σ.k_{a(j)} · p.δt
    return p
\end{verbatim}

各非零位置型軸 \(j\) に対して：

\[V[j].\theta_0 \mathrel{+}= \sigma.k_{a(j)} \cdot \delta t\]

これだけ。\texttt{σ}, \texttt{δV}, \texttt{δt} は \texttt{step(p)}
では一切更新されない。

\textbf{v2.0 からの変更}: v2.0 では \texttt{δt\ +=\ δδt}
を含めていたが、v2.1 で削除（§5.2 参照）。

\subsubsection{§6.2 型整合検査}\label{ux578bux6574ux5408ux691cux67fb}

\begin{verbatim}
σ.k_axis · δt : SpatialFreq × DeltaT 
              = Linv256 × L256 
              = Scalar256 [1]                ✓

V[j].θ₀ + (σ.k · δt) : CyclePhase + Scalar256
                     = CyclePhase [1] (mod P/2)  ✓
\end{verbatim}

\subsubsection{§6.2.1 5
原則の充足（v2.1）}\label{ux539fux5247ux306eux5145ux8db3v2.1}

\texttt{step(p)} は処理関数の 5 基本原則をすべて満たす：

\begin{enumerate}
\def\labelenumi{\arabic{enumi}.}
\tightlist
\item
  \textbf{null 透過性}: 関数の最初で明示的に処理 ✓
\item
  \textbf{in-place 更新}: V{[}j{]}.θ₀ を直接更新、新しい Particle 非生成
  ✓
\item
  \textbf{型整合性}: D1 §2.3 の規則に準拠（§6.2 参照） ✓
\item
  \textbf{INPUT/OUTPUT 入れ替え可能性}: 関数は \texttt{p}
  を読み書きする対称構造 ✓
\item
  \textbf{時間反転による可逆性}: δt の符号反転で完全に巻き戻る：
\end{enumerate}

\begin{verbatim}
1 回目（δt = +δt₀）: V → V + σ.k · δt₀
δt 反転（δt → -δt₀）
2 回目（δt = -δt₀）: V → V + σ.k · δt₀ + σ.k · (-δt₀) = V  ✓
\end{verbatim}

これは {[}W3{]} の情報保存要件を構造的に保証する。

\subsubsection{§6.3
偏光・複数軸位相差の表現}\label{ux504fux5149ux8907ux6570ux8ef8ux4f4dux76f8ux5deeux306eux8868ux73fe}

複数の非零軸を持つ粒子（例：偏光を持つ一方向光）：

\begin{verbatim}
σ_polarized_light = (+k_x, +k_y, 0, 0, +R_γ, 0)
V = [wave_x, wave_y]
\end{verbatim}

偏光位相差として：

\[\Delta\varphi = V[0].\theta_0 - V[1].\theta_0 \in \text{CyclePhase}\]

{\def\LTcaptype{none} % do not increment counter
\begin{longtable}[]{@{}ll@{}}
\toprule\noalign{}
偏光状態 & \(\Delta\varphi\) \\
\midrule\noalign{}
\endhead
\bottomrule\noalign{}
\endlastfoot
線偏光（同方向） & \(0\) \\
線偏光（反方向） & \(P/4 = 180°\) \\
円偏光（左） & \(+P/8 = +90°\) \\
円偏光（右） & \(-P/8 = -90°\) \\
楕円偏光 & 任意の整数値 \\
\end{longtable}
}

全て int256 整数で正確に表現される（\(15°\) の整数倍は \(P/48\)
の整数倍として誤差なし）。

\subsubsection{§6.4 時間の取り扱い：δt は処理ステップ、時間は
V{[}t{]}.θ₀}\label{ux6642ux9593ux306eux53d6ux308aux6271ux3044ux3b4t-ux306fux51e6ux7406ux30b9ux30c6ux30c3ux30d7ux6642ux9593ux306f-vt.ux3b8ux2080}

\textbf{重要原則}: δt
は時間ではなく、普遍的な処理ステップ。時間として観測されるのは粒子ごとの
V{[}t-軸{]}.θ₀ の累積値。

電子の例（σ.k\_t ≠ 0）:

\begin{verbatim}
σ_electron = (+k_x, 0, 0, +k_t, 0, 0)
V = [wave_x, wave_t]
1 ステップ:
  V[0].θ₀ += σ.k_x · δt   -- x 軸位相が進む（運動量相当）
  V[1].θ₀ += σ.k_t · δt   -- t 軸位相が進む（「電子が経験する時間」）
\end{verbatim}

光子の例（σ.k\_t = 0）:

\begin{verbatim}
σ_photon = (+k_x, +k_y, 0, 0, +R_γ, 0)
V = [wave_x, wave_y]      -- t 軸成分なし
-- σ.k_t = 0 ゆえ V に t 軸 OneDWave は存在しない
-- 結果: 光子の固有時間ゼロが構造的に保証される
\end{verbatim}

\subsubsection{§6.5
加速度効果は相互作用処理で扱う（v2.1）}\label{ux52a0ux901fux5ea6ux52b9ux679cux306fux76f8ux4e92ux4f5cux7528ux51e6ux7406ux3067ux6271ux3046v2.1}

v2.0 では「δδt は背景 R(r) や相互作用によって決まる」としていたが、v2.1
では：

\begin{itemize}
\tightlist
\item
  \textbf{加速度効果は単一粒子の基本演算では扱えない}（保存則・可逆性に反する、§5.2
  参照）
\item
  \textbf{すべての加速度効果は相互作用処理で実現}（粒子間または粒子-Mesh
  間で σ や δt が交換される）
\end{itemize}

具体的な相互作用規則は別稿 D2
で定義する。重力時間遅れ・偏光回転・周波数シフトは、すべて相互作用イベントの集積として保存則整合的に表現される。

\subsubsection{§6.6 必要としない演算（v1.0
と同じ）}\label{ux5fc5ux8981ux3068ux3057ux306aux3044ux6f14ux7b97v1.0-ux3068ux540cux3058}

以下の演算は本仕様の範囲では必要としない：

\begin{enumerate}
\def\labelenumi{\arabic{enumi}.}
\tightlist
\item
  \textbf{除算}: \(R/R_1\) は \(R \cdot R_1^{-1}\) として乗算に変換
\item
  \textbf{平方根}: 連続体的概念。離散系では整数値が直接的に物理量を決定
\item
  \textbf{三角関数}: 連続体 sine-Gordon の道具。離散変調表で代替可能
\item
  \textbf{浮動小数点演算}: 全フィールド
  int256、丸め誤差は構造的に発生しない
\end{enumerate}

cos/sin
が必要な箇所（例：縦波振動の変調）は、有限段階の離散変調表（lookup
table）で代替する。

\begin{center}\rule{0.5\linewidth}{0.5pt}\end{center}

\subsection{§7 Mesh 構造の新設（v2.0
で追加）}\label{mesh-ux69cbux9020ux306eux65b0ux8a2dv2.0-ux3067ux8ffdux52a0}

\subsubsection{§7.1
縦波モードの分離}\label{ux7e26ux6ce2ux30e2ux30fcux30c9ux306eux5206ux96e2}

{[}W8{]} 表 A
の解釈において、ヒッグス（\(\sigma_H = (0, 0, +k_z, 0, 0, 0)\)）は他のスピン1ボソンと異なり
\(\sigma.R = 0\) かつ \(\text{spin} = 0\)
である。これは構造的に「位置位相を振動させる横波粒子」ではなく、\textbf{「メッシュ自身の縦波振動」}
として扱うのが自然である。

同様に、グラビトン（重力波）も時空計量の縦波的振動として、Mesh
の振動モードと解釈できる。

これらの縦波モードを Particle
として扱うと、データ構造に無理が生じる（OneDWave
は本質的に位相回転を記述する型であり、軸方向の伸縮振動を直接表現できない）。本稿では
\textbf{Mesh 構造を新設し、縦波モードを Particle と分離管理} する。

\subsubsection{§7.2 Mesh
の構造（v2.1）}\label{mesh-ux306eux69cbux9020v2.1}

\begin{verbatim}
Mesh = (
    -- 背景スケール
    R_0           : L256,             -- 中心投影の中心長さ（背景）
    
    -- 単一軸縦波モード（spin 0、SO(3) 対称性破れ）
    H_x           : OneDWave,         -- x 軸縦波
    H_y           : OneDWave,         -- y 軸縦波
    H_z           : OneDWave,         -- z 軸縦波 = ヒッグス
    H_t           : OneDWave,         -- t 軸縦波
    
    -- 多軸縦波モード（spin 2、重力波の偏光）
    H_xy_anti     : OneDWave,         -- xy 反位相 = 重力波 + 偏光
    H_xy_quarter  : OneDWave,         -- xy 90° 位相 = 重力波 × 偏光
    H_yz_anti     : OneDWave,         -- yz 反位相
    H_xz_anti     : OneDWave,         -- xz 反位相
    
    -- 縦波モードの差分
    δH_x, δH_y, δH_z, δH_t            : OneDWave,
    δH_xy_anti, δH_xy_quarter         : OneDWave,
    δH_yz_anti, δH_xz_anti            : OneDWave,
    
    -- グローバル処理ステップ
    δt_mesh       : DeltaT,
)
\end{verbatim}

\textbf{v2.0 からの変更}:

\begin{itemize}
\tightlist
\item
  \texttt{H\_xyz\_inphase}（xyz 同位相）を削除し、軸別単一軸モード
  \texttt{H\_x,\ H\_y,\ H\_z,\ H\_t} に置換（理由：xyz
  等方解釈ではヒッグスの世代質量階層を説明できない、{[}W11{]} §8.5
  参照）
\item
  ヒッグスは \texttt{H\_z}（z 軸縦波）として位置付けられる
\item
  \texttt{δδt\_mesh} を削除（理由：Particle
  と同じく可逆性・保存則のため）
\end{itemize}

\subsubsection{§7.3
縦波モードの物理同定}\label{ux7e26ux6ce2ux30e2ux30fcux30c9ux306eux7269ux7406ux540cux5b9a}

{\def\LTcaptype{none} % do not increment counter
\begin{longtable}[]{@{}llll@{}}
\toprule\noalign{}
モード & 振動パターン & 物理同定 & spin \\
\midrule\noalign{}
\endhead
\bottomrule\noalign{}
\endlastfoot
H\_z & z 軸縦波（SO(3) 破れ） & \textbf{ヒッグス} & 0 \\
H\_x, H\_y, H\_t & 各軸縦波 & 物理同定未確定（候補） & 0 \\
H\_xy\_anti & x 伸び・y 縮みの反位相 & 重力波 + 偏光 & 2 \\
H\_xy\_quarter & xy 90° 位相回転 & 重力波 × 偏光 & 2 \\
H\_yz\_anti & yz 反位相 & 重力波（別偏光） & 2 \\
H\_xz\_anti & xz 反位相 & 重力波（別偏光） & 2 \\
\end{longtable}
}

\textbf{ヒッグス = z 軸縦波（H\_z）}: W8 表 A の
\(\sigma_H = (0, 0, +k_z, 0, 0, 0)\) における z 軸選択を SO(3)
対称性破れの方向として保持する。第 3 世代フェルミオンが z
軸を共有することで最強結合 → 最大質量を獲得する（{[}W11{]} §8 と整合）。

\textbf{重力波の 2 偏光 (+, ×)}: 多軸縦波モード（H\_xy\_anti,
H\_xy\_quarter, etc.）として表現。SO(3) 対称性は保たれる。

\subsubsection{§7.4 Mesh と Particle
の結合（v2.1）}\label{mesh-ux3068-particle-ux306eux7d50ux5408v2.1}

粒子の振動周波数 σ.k は Mesh
の状態によって実効的に変調される。質量生成・周波数シフトはこの変調を介した相互作用として表現される（詳細は別稿
D2 参照）。

実効周波数（質量生成の例）：

\[\sigma.k_{\text{effective}}(particle, mesh) = \sigma.k(particle) - \delta k_H(\text{Mesh.H\_z VEV})\]

z 軸を共有する第 3 世代フェルミオンが H\_z VEV
と最強結合し、最大質量を獲得する（{[}W11{]} §8 と整合）。

\subsubsection{§7.5 動力学（v2.1：基本演算は step
のみ）}\label{ux52d5ux529bux5b66v2.1ux57faux672cux6f14ux7b97ux306f-step-ux306eux307f}

Mesh の更新：

\begin{verbatim}
For each longitudinal mode H_*:
    H_*.θ₀ += (固有周波数) · δt_mesh
\end{verbatim}

Particle の更新（Mesh と独立、基本演算は step(p) のみ）：

\begin{verbatim}
For each particle p:
    step(p)   -- §6.1 の通り、V のみ更新
\end{verbatim}

Mesh-Particle 間の結合（質量生成・重力波など）は単一粒子の
\texttt{step(p)} ではなく、相互作用処理（D2 で定義）の中で σ や δt
を変化させる形で実現する。

\subsubsection{§7.6
メモリフットプリント（v2.1）}\label{ux30e1ux30e2ux30eaux30d5ux30c3ux30c8ux30d7ux30eaux30f3ux30c8v2.1-1}

\begin{verbatim}
R_0    : 256 bit
H_*    : 8 × 3 × 256 = 6,144 bit  -- 単一軸4 + 多軸4 = 8 モード
δH_*   : 8 × 3 × 256 = 6,144 bit
δt_mesh : 256 bit
合計   = 12,800 bit = 1,600 byte / 全宇宙
\end{verbatim}

宇宙全体で約 1.6 KB。Mesh モード数増加により v2.0 の 1 KB
から増えたが、粒子数 \(N\) に対して総メモリは \(608N + 1600\)
byte（粒子側は v2.0 から 32 byte 削減）。

\begin{center}\rule{0.5\linewidth}{0.5pt}\end{center}

\subsection{§8 W3
構造要件の充足確認（v2.1）}\label{w3-ux69cbux9020ux8981ux4ef6ux306eux5145ux8db3ux78baux8a8dv2.1}

{\def\LTcaptype{none} % do not increment counter
\begin{longtable}[]{@{}
  >{\raggedright\arraybackslash}p{(\linewidth - 4\tabcolsep) * \real{0.3333}}
  >{\raggedright\arraybackslash}p{(\linewidth - 4\tabcolsep) * \real{0.3333}}
  >{\raggedright\arraybackslash}p{(\linewidth - 4\tabcolsep) * \real{0.3333}}@{}}
\toprule\noalign{}
\begin{minipage}[b]{\linewidth}\raggedright
W3 要件
\end{minipage} & \begin{minipage}[b]{\linewidth}\raggedright
充足
\end{minipage} & \begin{minipage}[b]{\linewidth}\raggedright
根拠
\end{minipage} \\
\midrule\noalign{}
\endhead
\bottomrule\noalign{}
\endlastfoot
有界性 & ✓ & int256 で全位相量を表現可能（§1） \\
離散性 & ✓ & 全フィールドが整数型（§2--§5） \\
\textbf{情報保存（時間反転可逆性）} & \textbf{✓ (v2.1 で厳密化)} &
step(p) は δt 反転で完全に巻き戻る（§6.2.1） \\
ゼロ除算回避 & ✓ & 除算未使用（§6.6） \\
整数演算完結 & ✓ & 加算・乗算のみ、cos は離散変調表で代替（§7.4） \\
\textbf{6軸対称性} & \textbf{✓ (v2.0〜)} & 周波数解釈で全軸対等（§1.4,
§3） \\
\textbf{横波・縦波分離} & \textbf{✓ (v2.0〜)} & Particle vs
Mesh（§7） \\
\textbf{保存則の構造的維持} & \textbf{✓ (v2.1 で厳密化)} & step(p) は
σ・δt を変えない、保存則自明 \\
時間の非特権化 & ✓ & t は他軸と対等、時間は V{[}t{]}.θ₀（§6.4） \\
\textbf{5 原則すべて充足} & \textbf{✓ (v2.1)} & null
透過・in-place・型整合・対称性・可逆性（§6.2.1） \\
有限表現可能性 & ✓ & 粒子あたり ≤ 608 byte、全宇宙 +1.6 KB（§5.4,
§7.6） \\
\end{longtable}
}

\begin{center}\rule{0.5\linewidth}{0.5pt}\end{center}

\subsection{§9 v1.0
からの非互換性}\label{v1.0-ux304bux3089ux306eux975eux4e92ux63dbux6027}

\subsubsection{§9.1 解釈の変更}\label{ux89e3ux91c8ux306eux5909ux66f4}

\begin{itemize}
\tightlist
\item
  σ.k は「波数・加速度」から「振動周波数」へ統一
\item
  δt は「固有時間ステップ」から「処理ステップ周期」へ明示化
\item
  時間は V{[}t-軸{]}.θ₀ として現れる派生量
\end{itemize}

\subsubsection{§9.2
構造の拡張と整理}\label{ux69cbux9020ux306eux62e1ux5f35ux3068ux6574ux7406}

\textbf{v2.0 で追加}（v2.1 で継続）:

\begin{itemize}
\tightlist
\item
  Particle に δV を追加（v1.0 の Particle はサブセットとして含まれる）
\item
  Mesh 構造を新設（v1.0 にはなかった）
\end{itemize}

\textbf{v2.0 で追加し v2.1 で削除}:

\begin{itemize}
\tightlist
\item
  Particle の \texttt{δδt} フィールド（保存則・可逆性違反のため削除）
\item
  Mesh の \texttt{δδt\_mesh} フィールド（同上）
\item
  Mesh の \texttt{H\_xyz\_inphase}
  モード（観測される世代質量階層を説明できないため、\texttt{H\_z}
  等の軸別モードに置換）
\end{itemize}

\subsubsection{§9.3 W8 表 A
との整合}\label{w8-ux8868-a-ux3068ux306eux6574ux5408}

{[}W8{]} 表 A の 62 粒子分類は v2.1 でも有効。ヒッグスは \textbf{Mesh の
H\_z 励起}として位置付けられる（v2.0 で提案した H\_xyz\_inphase は v2.1
で訂正、{[}W11{]} §8.5 と整合）。標準模型 61 状態 + グラビトン 1 状態 =
62 という総数は変わらない。

\begin{center}\rule{0.5\linewidth}{0.5pt}\end{center}

\subsection{§10 結語}\label{ux7d50ux8a9e}

v2.1 の本質：

\begin{itemize}
\tightlist
\item
  \textbf{xyzt は周波数}（位相空間の自然な解釈）
\item
  \textbf{6 軸完全対称}（位置型 4 軸 + R + Q）
\item
  \textbf{横波（Particle）と縦波（Mesh）の分離}
\item
  \textbf{基本演算 step(p) は V のみ更新}（5
  原則すべて充足、加速度効果は相互作用処理に分離）
\item
  \textbf{時間は V{[}t{]}.θ₀}（位置位相として粒子ごとに進む）
\item
  \textbf{全演算は int256 整数で閉じる}
\end{itemize}

この構造で動力学規則（D2、別稿）を書けば：

\begin{itemize}
\tightlist
\item
  重力時間遅れ・偏光回転（粒子-Mesh 相互作用）
\item
  ヒッグス機構（Mesh.H\_z VEV と粒子の σ.k\_z 結合）
\item
  重力波放射（Mesh.H\_xy\_anti モード励起）
\item
  クーロン 2 体問題（粒子間相互作用）
\end{itemize}

が全て同じ枠組みで計算可能になる。すべての加速度効果は相互作用処理として保存則・可逆性を満たす形で実現される。

本稿は設計仕様であり、これらの動力学的応用は別稿に委ねる。

\begin{center}\rule{0.5\linewidth}{0.5pt}\end{center}

\subsection{参考文献}\label{ux53c2ux8003ux6587ux732e}

\begin{itemize}
\tightlist
\item
  {[}W3{]} Kihara, N. ``万物の理論が満たすべき構造要件の定式化.'' DOI:
  10.5281/zenodo.19601592 (2026).
\item
  {[}W8{]} Kihara, N. ``6次元超直方体の集合構造とその組合せ論的性質.''
  DOI: 10.5281/zenodo.19762134 (2026).
\item
  {[}W9{]} Kihara, N. ``sine-Gordon方程式 ── 位相的ソリトンの基礎理論.''
  DOI: 10.5281/zenodo.19650966 (2026).
\item
  {[}W10{]} Kihara, N. ``形不変波の4つのモード.'' DOI:
  10.5281/zenodo.19709798 (2026).
\item
  {[}W11{]} Kihara, N. ``形不変波の相互作用.'' DOI:
  10.5281/zenodo.19763463 (2026).
\item
  {[}M1{]} Kihara, N. ``中心投影の三層構造モデル: R--R₁--R₀.'' DOI:
  10.5281/zenodo.19691713 (2026).
\item
  {[}思考実験\_R軸{]} Kihara, N. ``6次元符号化 xyztRQ の再検討.'' DOI:
  10.5281/zenodo.19904714 (2026).
\end{itemize}

\begin{center}\rule{0.5\linewidth}{0.5pt}\end{center}

\textbf{著者}: 木原 範昭 (Noriaki Kihara) \textbf{ORCID}:
\href{https://orcid.org/0009-0004-6753-4020}{0009-0004-6753-4020}
\textbf{ライセンス}: CC BY 4.0

\end{document}
